% Rehearsal packet — team retrank, NTCIR-19 CAMEO CIR
% Source of truth for every number: docs/paper/ntcir19_cameo_retrank.tex (submitted, final).
\documentclass[11pt]{article}
\usepackage[a4paper,margin=2cm]{geometry}
\usepackage{booktabs,tabularx,longtable,amssymb}
\usepackage[shortlabels]{enumitem}
\usepackage{xcolor}
\usepackage[most]{tcolorbox}
\usepackage[hidelinks]{hyperref}
\usepackage{parskip}
\setlength{\emergencystretch}{2em}

\definecolor{accent}{RGB}{0,96,160}
\definecolor{warn}{RGB}{170,40,30}
\newtcolorbox{say}[1]{colback=accent!4,colframe=accent,fonttitle=\bfseries,
  title={#1},left=2mm,right=2mm,top=1mm,bottom=1mm,breakable}
\newtcolorbox{caution}{colback=warn!4,colframe=warn,left=2mm,right=2mm,top=1mm,bottom=1mm,breakable}
\setlist{itemsep=1pt,topsep=2pt}
\newcommand{\inpaper}{\textcolor{accent}{\small[paper]}}
\newcommand{\bg}{\textcolor{black!55}{\small[background --- don't quote]}}

\title{\textbf{Rehearsal Packet}\\[2pt]\large Team retrank --- NTCIR-19 CAMEO Composed Image Retrieval}
\author{Alessandro Magnani (Independent)}
\date{Prepared 24 November 2026}

\begin{document}
\maketitle
\tableofcontents
\bigskip

\begin{caution}
\textbf{Ground rule.} Every number you say aloud must be in the submitted paper.
Items marked \bg{} are project history for your own understanding: true, but not in
the paper, so do not present them as results. Items marked \inpaper{} are safe to quote.
\end{caution}

%=====================================================================
\section{Logistics}
\begin{itemize}
  \item \textbf{Talk slot:} 8 minutes total = \textbf{6 min talk + 2 min Q\&A}, hybrid CAMEO
        session (organizers' email of 27 Oct). Program:
        \url{https://research.nii.ac.jp/ntcir/ntcir-19/program20261016.html}.
  \item \textbf{Slides:} \texttt{docs/slides/retrank\_cameo\_slides.pdf} --- 7 slides, 16:9.
        Target \textbf{5:40--5:50} so you never get cut off.
  \item \textbf{Poster:} \texttt{docs/poster/NTCIR19-CAMEO-MagnaniA-poster.pdf} --- A0 portrait,
        fonts embedded. \textbf{Upload by 27 Nov (AoE)} via
        \url{https://forms.gle/nqv8jNeEYjB3nnMG8} (one poster per task: CAMEO and R2C2).
  \item \textbf{Bring:} the slides PDF on laptop + USB stick + emailed to yourself; a
        presenter timer; this packet.
\end{itemize}

%=====================================================================
\section{Competition setup at a glance \inpaper}
\begin{itemize}
  \item \textbf{Venue:} NTCIR-19, CAMEO task (Composed Access to Multimodal E-commerce
        Objects), \textbf{Composed Image Retrieval (CIR) subtask}. We cite the organizers'
        overview paper.
  \item \textbf{Task:} given a \textbf{reference product image} and a
        \textbf{natural-language modification} (a shopper request, e.g.\ ``the same product
        but cheaper'', ``a gel cleanser from brand X'', ``something for sensitive skin''), return
        a ranked list from a fixed catalogue.
  \item \textbf{Corpus:} \textbf{4,086} Vietnamese skincare products (cleansers, toners,
        sunscreens and more). It comes from the ViEcomRec dataset (originally 2,244 face
        cleansers), extended by the organizers. Each product has an image (some are short
        videos; we use the first frame), a brand, a numeric price in VND, and a translated
        name/description. The full collection has 387,034 reviews from 318,500 users.
  \item \textbf{Queries:} \textbf{20} official CIR topics: a reference product plus a
        Vietnamese request, with an English translation. There were \textbf{no labels before
        submission}.
  \item \textbf{Evaluation:} results are pooled across participants and judged by humans on
        a \textbf{3-point scale}; the metrics are \textbf{nDCG@5 / nDCG@10}. The judgments were
        released on 2026-08-17.
  \item \textbf{Submission:} up to \textbf{5 runs} per subtask, as JSON giving the ranked
        image filenames and product ids.
  \item \textbf{Participation:} \textbf{5 teams, 15 CIR runs}. We placed \textbf{3rd of 5}
        at both cutoffs (best run 0.740 / 0.714).
  \item \textbf{Our own evaluation setup:} the official set was small and unlabeled, so we
        built \textbf{Dev}: 429 templated queries with one exact answer each, across five rules
        (cheaper 20, pricier 20, format swap 150, brand swap 89, claim swap 150). Merged with
        LLM-paraphrased free-form queries, it gives 2,083. We also used an \textbf{Official-20
        proxy}: a cached Claude Sonnet vision judge. The three regimes are not comparable, and
        we \textbf{never tuned on the official queries}.
\end{itemize}

%=====================================================================
\section{Paper contributions \inpaper}
The paper's framing: ``a set of empirically grounded findings'' rather than a single new
model, built on a standard two-stage design (Stage-1 retrieval $\to$ Stage-2 VLM rerank).
\begin{enumerate}
  \item \textbf{Recall is the bottleneck.}
    \begin{itemize}
      \item The limit is whether the answer enters the Stage-1 pool at all.
      \item Dev R@50 plateaus at about 0.44 for visual Stage-1; with one gold answer, no
            reranker can exceed the pool's recall.
      \item Officially, three rerankers over the same pool span only 0.030 nDCG@10
            (per-topic ranges up to 0.22). Swapping Stage-1 with the reranker fixed gains
            +0.078 and +0.103.
    \end{itemize}
  \item \textbf{Structured metadata is a complementary recall lever.}
    \begin{itemize}
      \item An LLM attribute extractor plus symbolic brand, price, format and claim matching
            recovers products a visual encoder can't see.
      \item R@50 rises from 0.44 to 0.53 on structured query types.
      \item Price carries the gain: brand only 0.084 vs.\ brand + price 0.232.
      \item It is roughly neutral on the full Dev set, which is mostly free-form (0.115 vs.\
            0.119).
    \end{itemize}
  \item \textbf{Local open-weight VLM rerankers are competitive with a hosted API.}
    \begin{itemize}
      \item Three of the four local models beat Haiku on the proxy; Pixtral is marginally
            below.
      \item The marginal cost is near zero versus about \$11 per pass on the API, which moves
            the cost from per-token to flat GPU-hours. That helps both iteration speed and
            reproducibility.
    \end{itemize}
  \item \textbf{A caution: visual reranking may conflict with exact structural matches.} In
        one Dev comparison, reranking the metadata pool went from 0.200 to 0.173. The mechanism
        is unverified and needs a repeated controlled run.
  \item \textbf{Methodological:} an explicit \textbf{three-regime evaluation} (Dev / proxy /
        official) with stated validity threats, plus a documented account of what did not
        work.
\end{enumerate}

%=====================================================================
\section{The message in one breath}
\begin{say}{One sentence}
On CAMEO's skincare catalogue, getting the right products into the candidate pool mattered
more, in aggregate, than how we reranked them: product metadata, especially price, finds what
a visual encoder can't, and local open-weight rerankers are competitive with a hosted API at
near-zero marginal cost.
\end{say}

\textbf{Three claims} (in this order, everywhere):
\begin{enumerate}
  \item \textbf{Recall is the ceiling.} No reranker can beat the candidate pool's recall.
  \item \textbf{Price is the metadata signal that matters.} It is invisible to vision.
  \item \textbf{Local open-weight VLM rerankers are competitive} with a hosted API at
        $\sim$\$0 marginal cost.
\end{enumerate}

\textbf{Five numbers to know cold} \inpaper:
\begin{center}
\begin{tabularx}{\linewidth}{@{}lX@{}}
\toprule
0.740 / 0.714 & best official nDCG@5 / @10 (Run1); 3rd of 5 teams at both cutoffs \\
0.019 & spread of the top four teams on nDCG@10 (no ordering claimed among 2nd--4th) \\
0.436 $\to$ 0.527 & Dev R@50: best visual vs.\ metadata-fused Stage-1 \\
0.084 $\to$ 0.232 & Dev metadata-only nDCG@10: brand only vs.\ brand + price \\
0.030 vs.\ +0.078 / +0.103 & official: reranker swap vs.\ Stage-1 swap (nDCG@10) \\
\bottomrule
\end{tabularx}
\end{center}

%=====================================================================
\section{The 6-minute talk script}
About 125 words per minute. The times are cumulative targets. Read it aloud with a timer
at least three times; the bracketed lines are what to cut if you are running long.

\begin{say}{Slide 1 --- Title \hfill 0:00--0:20}
Hi, I'm Alessandro Magnani, team retrank. This talk is about composed image retrieval on
CAMEO's skincare catalogue, and the main message fits in one line: on this task, getting the
right products into the candidate pool mattered more, on average, than how we reranked them.
\end{say}

\begin{say}{Slide 2 --- The task \hfill 0:20--1:10}
You get a photo of a product and a shopper's request --- ``the same product but cheaper'',
``a gel cleanser from brand X'' --- and you rank 4,086 Vietnamese skincare products. There
are 20 official topics, judged by pooling all teams' runs, scored with nDCG at 5 and 10.

What makes it interesting is on the right: many requests are about price, brand or
format. Those live in the product's metadata, not in its pixels, so a visual encoder can't
see them.

One note before any numbers: I'll quote three kinds. Our own synthetic dev set, with one
exact answer per query; a Claude vision-judge proxy on the official queries; and the
official human scores. They are not comparable to each other.
\end{say}

\begin{say}{Slide 3 --- Approach \hfill 1:10--2:10}
Two stages. Stage 1 has two branches. The \emph{visual} branch is training-free: we
combine CLIP ViT-bigG embeddings of the reference image and the request, and each product's
image is fused with its name.

The \emph{metadata} branch uses an LLM to read the request and pull out the target brand,
price, format and claim. We then score products of that brand by how close their price is,
plus keyword matches --- and if no brand is confidently named, the branch stays silent.

We put confident metadata hits first, then the visual list. Stage 2 is an open-weight VLM,
Qwen2.5-VL-72B, served locally, which reorders the top twenty.
\end{say}

\begin{say}{Slide 4 --- Finding 1: recall is the ceiling \hfill 2:10--3:00}
First finding. A reranker can only reorder what Stage 1 found. With one gold answer per dev
query, the best any reranker can do is the pool's recall.

Look at recall at 50. The textual-inversion baseline, SEARLE, gets 0.15. The best visual
backbone gets 0.44 --- and visual backbones plateau there. Adding metadata takes it to
0.53; it's the only thing that broke the plateau.

So: raise the ceiling in Stage 1. Reranking only closes the gap beneath it.
\end{say}

\begin{say}{Slide 5 --- Finding 2: price; local rerankers \hfill 3:00--4:00}
Which metadata matters? Turning on one term at a time: brand alone gives 0.084; adding
price jumps to 0.232; adding the format and claim keywords on top actually lowers it a
little. Price is the signal --- it's exactly what ``cheaper'' asks for, and vision can't see
it. The honest limit: on our larger, mostly free-form dev set, metadata is roughly neutral.

On the right, rerankers. Same candidate pool, proxy judge: three of four local open models
score above the hosted Haiku API, Pixtral just below. The gaps are within noise, so I'd call
them competitive --- [but the cost changes from about eleven dollars a pass to roughly zero
once a GPU is rented; our whole GPU bill was about sixty dollars.]
\end{say}

\begin{say}{Slide 6 --- Official results \hfill 4:00--5:10}
Officially, our primary run --- metadata-fused plus Qwen-72B --- scored 0.740 and 0.714:
third of five teams at both cutoffs. At nDCG at 10 the top four are within 0.019, so I
wouldn't read an order into second to fourth.

Our five runs form a useful controlled comparison. Runs 2 to 4 share an identical candidate
pool with three different rerankers: they differ by only 0.03, not significant. Keep the
reranker fixed and change the Stage-1 pool, and you gain 0.08 and 0.10. So, in aggregate,
candidate generation is the bigger lever.

Two caveats. Per topic, the reranker choice still swings up to 0.22. And these contrasts
change the whole Stage 1 --- Run 5's brand-union added no candidates on the official topics ---
so the metadata-specific evidence comes from the dev ablation.
\end{say}

\begin{say}{Slide 7 --- Takeaways \hfill 5:10--5:50}
To wrap up. Recall, not reranking, caps this task in aggregate. Structured metadata ---
especially price --- recovers what vision can't see. And local open rerankers are competitive
at near-zero marginal cost. One caution: in one dev comparison, visually reranking the metadata
pool lowered quality from 0.200 to 0.173, and we haven't verified why. Next: better brand
resolution, repeating that audit, and validating the extractor on natural queries. Thank you
--- happy to take questions.
\end{say}

\textbf{If you're behind:} at 4:00 you should be starting Slide 6. If you're past 4:20,
skip the bracketed cost aside on Slide 5 and shorten the caveats to one sentence:
``these swap the whole Stage 1, so metadata's own effect is shown on dev.''

%=====================================================================
\section{Poster pitches}
Visitors arrive at different depths. Pick the version that matches their interest; point
at the blocks as you go.

\begin{say}{30 seconds (the hallway version)}
We did composed image retrieval on a Vietnamese skincare catalogue. The main lesson: getting
the right products into the candidate pool matters more than reranking. Here
[\emph{Finding 1}] a perfect reranker can't beat the pool's recall; product metadata,
mostly price, is what broke the visual ceiling [\emph{Finding 2}]; and local open rerankers
are competitive with a hosted API at near-zero marginal cost [\emph{Finding 3}]. Officially
we were third of five; at nDCG@10 the top four are within 0.019.
\end{say}

\begin{say}{2 minutes (the standard tour)}
\begin{enumerate}[leftmargin=*]
  \item \emph{Top row, The Task:} photo + shopper request $\to$ rank 4,086 skincare
        products; the catch is that many requests are about price or brand.
  \item \emph{Top row, Approach:} two stages --- a training-free CLIP branch plus a metadata branch
        where an LLM reads brand, price, format and claim from the request; metadata hits go
        first; a local VLM reranks the top 20.
  \item \emph{Evaluation Regimes:} three different scoreboards --- our synthetic dev set, an LLM-judge
        proxy, and the official human scores --- so don't compare across them.
  \item \emph{Finding 1:} recall at 50 goes 0.15 $\to$ 0.44 $\to$ 0.53; the visual
        backbones plateau, and only metadata breaks it.
  \item \emph{Finding 2:} brand alone 0.084, add price 0.232; the keyword bonuses add noise.
  \item \emph{Finding 3:} three of four local rerankers score above the API, within noise;
        Pixtral just below; near-zero marginal cost per pass.
  \item \emph{Official Results:} 0.740 / 0.714, third of five. Same pool with three rerankers
        moves 0.03; a new Stage-1 pool with the same reranker moves 0.08--0.10. Whole-Stage-1
        swaps, so metadata's own effect comes from Dev.
  \item \emph{Where It Still Fails:} brand swap; price-band topics.
\end{enumerate}
\end{say}

\begin{say}{5 minutes (for someone who works on CIR)}
The 2-minute tour, plus: the oracle argument (one gold answer $\Rightarrow$ oracle nDCG@10
$=$ recall at the reranked depth); the extractor's accuracy on templated queries
(brand 428/429, price 40/40); gold vs.\ extracted fields (0.211 vs.\ 0.131, with no loss
on the price rules); the Friedman test (p$\approx$0.52); the per-topic reranker range
(0.22); the caution that visually reranking the metadata pool lowered Dev quality
(0.200 $\to$ 0.173) with an unverified mechanism; and the negative results.
\end{say}

%=====================================================================
\section{Everything we did, and why}
In chronological order. Each item gives what we did, why, the result, and whether it's in
the paper.

\subsection*{Phase 1 --- Framing (May 2026)}
\begin{enumerate}[leftmargin=*]
  \item \textbf{Read the organizers' reading list} (MagicLens, Pic2Word/SEARLE, UNION, CASE,
        CIRCO). \emph{Why:} to know the zero-shot CIR state of the art before picking an approach.
        \emph{Decision:} zero-shot, two-stage (retrieve $\to$ rerank), with a training-free visual composition
        that requires no CIR training.
  \item \textbf{Translate the Vietnamese requests to English} (VN$\to$EN) and use English CLIP.
        \emph{Why:} cheapest path, and the queries come with English translations.
        \inpaper
\end{enumerate}

\subsection*{Phase 2 --- Baseline and measurement (May--June)}
\begin{enumerate}[resume,leftmargin=*]
  \item \textbf{SEARLE textual-inversion baseline} (CLIP ViT-L/14). \emph{Why:} built for
        zero-shot CIR, with light compute. \emph{Result:} Dev R@20 0.084, R@50 0.154. It
        became the weak reference pool. \inpaper
  \item \textbf{Built our own development set.} \emph{Why:} the 20 official queries had no
        labels, so we couldn't measure recall. 429 templated queries over five rules
        (cheaper 20, pricier 20, format swap 150, brand swap 89, claim swap 150), each with
        one answer known by construction, plus LLM paraphrases (Claude Sonnet,
        $\approx$\$40) for 2,083 queries in total, with two held-out splits. \inpaper
  \item \textbf{A proxy judge for the official queries} (cached Claude Sonnet vision judge,
        graded). \emph{Why:} the only quality signal on the real queries before results came
        out. \emph{Known weakness:} lenient, uncalibrated, n=20, so we used it only for
        direction. \inpaper
  \item \textbf{First hosted-API reranking} (Sonnet/Haiku listwise). \emph{Result:} it helped
        the proxy judge but did nothing on Dev, the first sign that reranking can't recover
        answers missing from the pool. \bg
\end{enumerate}

\subsection*{Phase 3 --- Attacking recall with better encoders (June)}
\begin{enumerate}[resume,leftmargin=*]
  \item \textbf{Stage-1 backbone sweep:} SEARLE-OTI (R@20 0.142), ViT-H/14 (0.343),
        ViT-bigG-14 + name fusion (0.359; R@50 0.436), EVA02-L (0.294). \emph{Why:} to test
        whether recall was limited by the encoder. \emph{Result:} about 4$\times$ better than
        SEARLE, then a plateau near 0.44 R@50. The visual backbones we evaluated all plateaued there.
        \inpaper
  \item \textbf{Name fusion:} blend each product's image vector with its product-name text
        vector ($a{=}0.2$). \emph{Why:} lets brand or format words in the request match the
        name. \emph{Result:} small help at low weight; high weight hurt. \inpaper
  \item \textbf{Moved reranking to local open-weight VLMs} (vLLM on rented A100s).
        \emph{Why:} API reranking cost $\sim$\$11 per pass, was rate-limited, and scaled with
        every experiment. \emph{Result:} Qwen-72B 0.913, InternVL3-38B 0.907, Qwen-32B 0.894,
        Haiku 0.872, Pixtral 0.869 on the proxy. GLM-4.5V failed to load. \inpaper
  \item \textbf{Symbolic brand-union for brand-swap queries.} \emph{Why:} the purely visual pool
        rarely contained the different-brand target.
        \emph{Result:} brand-swap recall $\sim$7.5$\times$; putting brand into the embedding
        instead was worse. \inpaper
  \item \textbf{Automated configuration search} (an LLM proposes, a GPU runs, 21 configs).
        \emph{Why:} cheap local reranking made each trial nearly free. \emph{Result:} best
        was bigG + brand-union + local VLM at top-20; it stopped improving, so we needed a
        new lever, not more tuning. \inpaper
  \item \textbf{Dead ends:} RRF ensembling of backbones, rewriting queries into richer
        captions, using format/claim as hard filters. All gave no gain, or collapsed the pool.
        \inpaper
\end{enumerate}

\subsection*{Phase 4 --- The new lever: metadata (June)}
\begin{enumerate}[resume,leftmargin=*]
  \item \textbf{LLM attribute extractor} (Claude Haiku, text-only). \emph{Why:} the corpus
        has brand, numeric price and descriptions that vision ignores; the request has to be
        turned into structured targets first. \emph{Result:} brand 428/429, price 40/40 on
        the queries that state one, format/claim produced on 150/150 (coverage, not
        correctness). \inpaper
  \item \textbf{Symbolic scoring + metadata-first fusion.} Brand gate; score = 1 + 2 per
        format hit + 2 per claim hit + Gaussian price proximity ($\sigma{=}0.4$, relative to
        the target price); abstain without a confident brand; metadata list first, then
        visual. \emph{Why:} a brand/price match is more trustworthy than a visual neighbour,
        and the two score scales aren't comparable, so we order rather than blend.
        \emph{Result:} Dev R@50 0.527, nDCG@10 0.200 (vs.\ 0.436 / 0.157 visual). \inpaper
  \item \textbf{Rerank-on-metadata caution.} Visually reranking the metadata pool lowered Dev
        exact-match (0.200 $\to$ 0.173) while the same reranker helps a visual pool (gold
        item moved up 24 times, down 13). The reranked metadata run wasn't saved, so the
        mechanism is unverified. \emph{Decision:} Run1 still used the reranker, because the
        graded proxy (closer to human judging) favoured it (0.923 on the proxy). \inpaper
  \item \textbf{MagicLens attempt.} Adapter written and weights obtained, but blocked by a
        pinned mid-2024 JAX/scenic environment we couldn't rebuild. Not a measured result.
        \inpaper
\end{enumerate}

\subsection*{Phase 5 --- Submission (July 10)}
\begin{enumerate}[resume,leftmargin=*]
  \item \textbf{Five runs, chosen for diversity} across Stage-1 pools and rerankers, to hedge
        against an unknown test distribution. Run1 (metadata-fused + Qwen-72B) was primary.
        The organizers confirmed no resubmission, so these runs are final. \emph{Honest
        note:} the proxy judge clustered all strong configs at $\sim$0.88--0.91, so it
        couldn't tell Stage-1 pools apart. That's why three SEARLE-pool runs made the
        cut. \bg
\end{enumerate}

\subsection*{Phase 6 --- Official analysis and paper (Aug--Oct)}
\begin{enumerate}[resume,leftmargin=*]
  \item \textbf{Verified each run's pipeline} by matching submitted candidate sets to the
        archived pools (Jaccard 1.00; SEARLE vs.\ bigG pools overlap only 0.22). \emph{Why:}
        our own run labels had to be confirmed before we drew any conclusions. \inpaper
  \item \textbf{Controlled comparisons from the official runs:} reranker-only (Runs 2--4,
        spread 0.030, Friedman p$\approx$0.52, per-topic up to 0.22) vs.\ matched-reranker
        Stage-1 swaps (+0.078, +0.103; 13/2/5 and 15/2/3 win/tie/loss). \inpaper
  \item \textbf{Price-band topics:} 11 topics at 0.586 vs.\ 9 at 0.870. Their recall
        isn't lower (0.512 vs.\ 0.492), so the gap is ordering, not coverage. Two of three
        comparable teams show the same pattern, so it's suggestive, not a task property.
        \inpaper
  \item \textbf{Paper analyses:} component ablation, gold vs.\ extracted fields, rank-shift
        audit, oracle bound. Reviewer fixes for the camera-ready (corpus 4,086 skincare; no
        platform name). \inpaper
\end{enumerate}

\textbf{Total spend} \inpaper: $\sim$\$235 hosted-LLM API (early reranking, dev-set
generation, proxy judge, extractor) + $\sim$\$60 GPU rental.

%=====================================================================
\section{Q\&A bank}
Keep answers to 20--30 seconds. With only 2 minutes of Q\&A you'll get 2--3 questions.

\begin{longtable}{@{}p{0.33\linewidth}p{0.63\linewidth}@{}}
\toprule
\textbf{Question} & \textbf{Answer} \\ \midrule \endhead
Does metadata explain your official gain? &
Not by itself. On the official set, our Stage-1 contrasts change the whole retriever
family, SEARLE to bigG-based, so they don't isolate metadata. The metadata attribution
comes from the controlled ablation on dev. \\ \addlinespace
Is the ordering among teams 2--4 meaningful? &
It isn't established. At nDCG@10 the top four are within 0.019 on 20 topics, and we ran no
significance or equivalence tests, so we don't infer an order. (At nDCG@5 the spread is
larger.) \\ \addlinespace
Why three runs on the weak SEARLE pool? &
We chose the five runs for diversity, and before results our proxy judge couldn't separate
the Stage-1 pools. It turned out to be useful: those three runs are a clean
reranker-only comparison on official data. \\ \addlinespace
Reranking hurt the metadata pool on dev --- why rerank Run1? &
That was one recorded dev comparison, and the mechanism isn't verified because the reranked
run wasn't saved. Dev is exact single-answer; the graded proxy favoured reranking. We report
it as a caution and an experiment to run. \\ \addlinespace
Your dev set is synthetic --- isn't that biased? &
Yes, it reflects the generator's notion of the right answer. That's why we use it only as a
relative recall signal, check direction on the proxy, and never tuned on the official
queries. \\ \addlinespace
Why training-free, not a trained CIR model? &
The training-free composition needs no CIR training, and the evidence pointed at recall from
metadata, not a better visual embedding. MagicLens, the purpose-built option, was blocked
by its environment. \\ \addlinespace
How accurate is the LLM extractor on real queries? &
On templated dev queries: brand 428/429, price 40/40. We haven't measured it on the
natural official queries; that's future work. \\ \addlinespace
What about ``same brand'' when the brand isn't named? &
The extractor is text-only, so it resolves the brand only if the text names it. A lookup
of the reference product's brand would handle the rest; we left that to the fusion side. \\ \addlinespace
Why a Gaussian on price? &
It peaks at the target price and decays with the relative difference, so ``around 49k''
prefers products near 49k. It's a fixed bandwidth, $\sigma=0.4$. \\ \addlinespace
Why listwise reranking? &
One call ranks all candidates jointly, so the model compares them against the request, and
it costs one generation per query. \\ \addlinespace
Why is the proxy so much higher than the official score? &
Different ground truth: a lenient graded LLM judge over our own runs vs.\ pooled human
judgments. The regimes aren't comparable. \\ \addlinespace
Why are price-band topics hard if recall is fine? &
Many products plausibly fit; the difficulty is putting the best ones at the top. It's
descriptive. Those topics also combine brand, format and ingredient constraints. \\ \addlinespace
The 0.22 per-topic reranker range? &
The 0.03 aggregate hides cancellation: individual topics are quite sensitive to the
reranker even though the averages match. \\ \addlinespace
Cost? &
About \$235 in hosted-API calls overall and about \$60 of GPU. Once a GPU is rented, an
extra rerank pass is essentially free, vs.\ about \$11 per pass on the API. \\ \addlinespace
How does this relate to UNION (the organizers' work)? &
UNION learns a lightweight target representation that unifies CIR and sketch retrieval. We
asked a different question: how far non-visual product metadata carries retrieval on a
structured catalogue. \\ \addlinespace
What would you do next? &
Better brand resolution (brand swap is the hardest rule); a reranker that keeps confident
metadata matches; one more rerank pass to audit the metadata demotion item by item;
validating the extractor on natural queries. \\
\bottomrule
\end{longtable}

%=====================================================================
\section{Don't say}
\begin{caution}
\begin{itemize}
  \item ``We're effectively second'' or any ordering among teams 2--4; the 0.019 figure
        is nDCG@10 only.
  \item ``Metadata gave +0.078 / +0.103 officially.'' Those are whole-Stage-1 swaps, and
        Run5's brand-union added no official candidates.
  \item ``Local models are better than the API.'' Say \emph{competitive}; the gaps are within noise.
  \item ``Reranking hurts metadata'' as a fact. It's a caution with an unverified mechanism.
  \item ``Price is hard for every team.'' It holds for two of three comparable teams.
  \item ``Metadata helps everywhere.'' On the full 2,083-query Dev set (79\% free-form) it's
        roughly neutral (0.115 vs.\ 0.119). We have no free-form-only number.
  \item The old corpus description (``4,075 face cleansers'') or the name of the source
        platform. Say \textbf{4,086 skincare products from a Vietnamese e-commerce platform}.
  \item Internal-only numbers from project notes that aren't in the paper.
\end{itemize}
\end{caution}

%=====================================================================
\section{Rehearsal checklist}
\begin{itemize}
  \item[$\square$] Read the script aloud with a timer 3$\times$; land at 5:40--5:50.
  \item[$\square$] Practise the 30-second poster pitch until you don't need notes.
  \item[$\square$] Say the five numbers from memory.
  \item[$\square$] Practise five Q\&A answers out loud, at 20--30 seconds each.
  \item[$\square$] Hybrid check: screen-share the slides PDF, mic, presenter timer.
  \item[$\square$] Upload the CAMEO poster (and R2C2) by 27 Nov AoE.
\end{itemize}

\end{document}
